% ============================================
% PREÂMBULO COMUM DOS SLIDES DE FM
% Usado por FM/main_fm.tex e pelo main.tex combinado do Módulo 2
% ============================================
\usepackage{pgfplots}
\pgfplotsset{compat=1.18}
\usetikzlibrary{decorations.pathmorphing, decorations.pathreplacing, arrows.meta,
                positioning, calc, fit, backgrounds, shapes.geometric, tikzmark}

% ============================================
% RECURSOS PARA DERIVAÇÃO ESTILO QUADRO
% (setas de anotação + destaques + overlays)
% ============================================
\colorlet{cRed}{red!65!black}
\colorlet{cBlue}{blue!65!black}
\colorlet{cGreen}{green!45!black}
\colorlet{cOrange}{orange!85!black}
\colorlet{cPurple}{purple!70!black}

\tikzset{
  nota/.style={draw=#1, text=#1, fill=white, rounded corners=2pt, inner sep=3pt,
               font=\scriptsize, align=center, line width=0.6pt},
  nota/.default=cRed,
  setinha/.style={-{Stealth[length=2.2mm,width=1.6mm]}, line width=0.8pt,
                  shorten >=2pt, shorten <=2pt},
}

% Destaque de um trecho de equação. Depende da baseline, então fica inline.
\providecommand{\hl}{}
\renewcommand{\hl}[2][cRed]{%
  \tikz[baseline=(HLX.base)]{%
    \node[inner sep=1.5pt, outer sep=0pt, fill=#1!18,
          rounded corners=1.5pt] (HLX) {$\displaystyle #2$};}}
% Setas de anotação sobre as equações, no estilo quadro.
% Uso: \anotacoes{ \only<2->{ \node[nota=cRed] ... ; \draw[setinha,cRed] ... ; } }
\providecommand{\anotacoes}[1]{%
  \begin{tikzpicture}[remember picture, overlay]#1\end{tikzpicture}}

% Espaço menor em volta das equações exibidas, para frames densos.
% Usar depois de \small, que redefine esses comprimentos.
\newcommand{\compacto}{%
  \setlength\abovedisplayskip{4pt}\setlength\belowdisplayskip{4pt}%
  \setlength\abovedisplayshortskip{2pt}\setlength\belowdisplayshortskip{2pt}}
